% Run GCT_KRON_20261120; GPT-6 Astra Pro; effort not exposed.
\section{An exact charge-block reduction of the permanent residue Jacobian}

\subsection*{QA033-CB: a monomial determinant identity and its remaining gate}
Work over a field of characteristic zero, let \(m\geq3\), and use
the distinct parameters, first-jet kernel \(K\), and inverse-factorial
binary coefficient space from (RJ.1)--(RJ.10).
This section proves an exact reduction. It does not assert that its
last square map is invertible for arbitrary \(m\).
Put
\[
 \begin{gathered}
 k=m-1,\qquad M=\binom m2,\qquad
 H=\mathbb C[t]_{\leq m-2},\\
 P=\langle p_1,\ldots,p_m\rangle,\qquad
 Q=\langle q_1,\ldots,q_{m-1}\rangle,\qquad q_m=0.
 \end{gathered}
                                                               \tag{CB.1}
\]
Here \(P,Q\) are spaces of transverse linear forms, so the
target quadratic coefficient space is
\(H\otimes\operatorname{Sym}^2(P\oplus Q)\).
Let \(z\in P\), with the tested choice \(z=p_m\).
For a scalar parameter \(\lambda\), define the off-diagonal
entries of \(B_\lambda\in K\otimes(P\oplus Q)\) by
\[
 (B_\lambda)_{ij}=p_i+p_j+q_i-q_j
                         +\lambda\operatorname{sgn}(j-i)z.
                                                               \tag{CB.2}
\]
Complete the diagonal by (RJ.7). The added skew terms require
no diagonal correction.

Let \(E_S,E_A\) be copies of the unordered-edge space, both
of dimension \(M\), represented in \(K\) by
\[
 \kappa^S_{ab}=\tfrac12(K_{ab}+K_{ba}),\qquad
 \kappa^A_{ab}=\tfrac12(K_{ab}-K_{ba})\quad(a<b).
                                                               \tag{CB.3}
\]
They give the direct sum \(K=E_S\oplus E_A\).
Write the three gradient maps supplied by (O-031.3) as
\[
 \begin{aligned}
 s&:E_S\longrightarrow H\otimes P,\\
 a&:E_A\longrightarrow H\otimes Q,\\
 c&:E_A\longrightarrow H\otimes P.
 \end{aligned}                                                   \tag{CB.4}
\]
Explicitly, with \(A=a_a,B=a_b,D=B-A\),
\[
 \begin{aligned}
 s(\kappa^S_{ab})
   &=-\frac{A+B}{D}(p_a-p_b)P_{ab}
           -\frac{B e_a-A e_b}{D}\Delta_{ab},\\
 a(\kappa^A_{ab})&=-(q_a-q_b)P_{ab},\\
 c(\kappa^A_{ab})&=-zQ_{ab}.
 \end{aligned}                                                   \tag{CB.5}
\]
Thus the quadratic differential \(\mathcal J_\lambda\) sends
\(e\otimes\ell\) to \(s(e)\ell\) in the symmetric edge sector
and to \((a(e)+\lambda c(e))\ell\) in the skew edge sector.
This is the differential of the inverse-factorial quadratic
coefficient of the actual permanent. Passing to residues and
to the normalization in (RJ.13) changes bases and a fixed
nonzero scalar only, independently of \(\lambda\).

\paragraph{The unchanged pure-\(Q^2\) isomorphism.}
The multiplication map
\[
 D:E_A\otimes Q\longrightarrow H\otimes\operatorname{Sym}^2Q,
                         \qquad e\otimes\ell\longmapsto a(e)\ell
                                                               \tag{CB.6}
\]
is an isomorphism. This is the uniform pure-skew assertion of
O-031. For clarity, its full argument is as follows.
Apply \((1+t\partial_t)^{-1}\) to the binary coefficients,
then divide by \(\Delta\) and take residues. Up to an individual
nonzero factor for each edge, the resulting map is
\[
 e_{ab}\otimes\ell\longmapsto
              (\epsilon_a-\epsilon_b)\otimes(q_a-q_b)\ell.
                                                               \tag{CB.7}
\]
The residue space is \(\{(v_i):\sum_i v_i=0\}\).
The dimensions of source and target are both \(Mk\).
A dual vector consists of symmetric bilinear forms \(T_i\)
on the space with basis \(q_1,\ldots,q_{m-1}\), modulo a common
form. Set \(T_m=0\). The edge equations are
\((T_a-T_b)(q_a-q_b,\xi)=0\) for all \(\xi\).
Edges \(\{i,m\}\) imply \(T_i(q_i,\xi)=0\), and the remaining
edges imply \(T_i(q_j,\xi)=-T_j(q_i,\xi)\).
Putting \(\xi=q_j\) kills the other diagonal coefficients.
For three distinct indices, symmetry and three successive
edge identities give
\[
 T_i(q_j,q_l)=-T_j(q_i,q_l)
             =T_l(q_i,q_j)=-T_i(q_j,q_l).
\]
All coefficients vanish because \(2\) is invertible. Dual
injectivity and equal dimensions prove (CB.6).

\paragraph{An unconditional square reduced map.}
Define the source and target spaces
\[
 \begin{aligned}
 G&=(E_S\otimes Q)\oplus(E_A\otimes P),&
 U&=E_S\otimes P,\\
 W&=H\otimes P\otimes Q,&
 Y&=H\otimes\operatorname{Sym}^2P.
 \end{aligned}                                                   \tag{CB.8}
\]
The relevant maps are
\[
 \begin{aligned}
 A:G&\longrightarrow W,&
 A(u,v)&=s(u)+a(v),\\
 F:U&\longrightarrow Y,& F(u)&=s(u),\\
 C:G&\longrightarrow Y,& C(u,v)&=c(v).
 \end{aligned}                                                   \tag{CB.9}
\]
In these formulas the unused tensor factor is multiplied into
the indicated gradient, as in (CB.6). In particular \(C\)
vanishes on \(E_S\otimes Q\).

Order source blocks as \((E_A\otimes Q),G,U\), and target
blocks as \(H\otimes\operatorname{Sym}^2Q,W,Y\).
There is a further fixed map \(E:E_A\otimes Q\to W\), obtained
by multiplying \(c(e)\) by the \(Q\)-form, and the full matrix is
\[
 \mathcal J_\lambda=
 \begin{pmatrix}
       D&0&0\\
       \lambda E&A&0\\
       0&\lambda C&F
 \end{pmatrix}.
                                                               \tag{CB.10}
\]
This follows directly from the degrees in \(P\) and \(Q\):
the symmetric gradient has degree \((1,0)\), the charge-zero
skew gradient has degree \((0,1)\), and its charge term has
degree \((1,0)\).
Subtracting \(\lambda ED^{-1}\) times the first row block from
the second eliminates its first block without altering any
other block. Therefore
\[
 \det\mathcal J_\lambda=(\det D)\det R_\lambda,\qquad
 R_\lambda=\begin{pmatrix}A&0\\ \lambda C&F\end{pmatrix}
                      :G\oplus U\longrightarrow W\oplus Y.
                                                               \tag{CB.11}
\]
The reduced matrix is square. Its dimensions are
\[
 \begin{array}{c|c}
 \dim G=M(2m-1)&\dim W=m(m-1)^2\\
 \dim U=mM&\dim Y=(m-1)m(m+1)/2,
 \end{array}
                                                               \tag{CB.12}
\]
In particular,
\(\dim G-\dim W=\dim Y-\dim U=M\).
No rank assertion about \(A\) or \(F\) is used in
(CB.10)--(CB.12).

\paragraph{The exact determinant identity.}
Over the rational-function field in \(\lambda\), multiply
the \(W\)-rows of \(R_\lambda\) by \(\lambda\) and the
\(G\)-columns by \(\lambda^{-1}\). The resulting matrix is
exactly \(R_1\). The product of the determinant factors is
\(\lambda^{\dim W-\dim G}=\lambda^{-M}\). Consequently
\[
 \boxed{\det\mathcal J_\lambda
                 =\lambda^M\det\mathcal J_1.}
                                                               \tag{CB.13}
\]
Both sides are polynomials, so this also holds at \(\lambda=0\).
The identity holds with either side identically zero and for
every fixed choice of \(z\in P\); in particular it makes no
implicit nonvanishing assumption. Every nonzero charge scalar
works if one of them works.

\paragraph{Precisely when the last gate has dimension \(M\).}
The following three assertions together are equivalent to
invertibility of \(R_1\), and hence of \(\mathcal J_1\):
\[
 \begin{gathered}
 A\text{ is surjective},\qquad F\text{ is injective},\\
 \overline C:\ker A\longrightarrow\operatorname{coker}F,
                \qquad v\longmapsto C(v)\bmod\operatorname{im}F,
       \quad\text{is an isomorphism}.
 \end{gathered}                                                \tag{CB.14}
\]
Under the first two assertions, both spaces in the last line
have dimension \(M\). To verify the equivalence, an element
\((v,u)\) is in \(\ker R_1\) exactly when
\(Av=0\) and \(Cv+Fu=0\). If (CB.14) holds, these equations
force first \(v=0\), then \(u=0\). Equal dimensions give
invertibility. Conversely, invertibility forces \(A\) to
surject onto the first target block and forces \(F\) to be
injective. An element in \(\ker\overline C\) then lifts to
an element of \(\ker R_1\), and is zero. Equal dimensions
make \(\overline C\) an isomorphism.

The charge-zero nullity has the exact formula
\[
 \dim\ker\mathcal J_0
   =M+(\dim W-\operatorname{rank}A)
                         +(\dim U-\operatorname{rank}F).
                                                               \tag{CB.15}
\]
Indeed \(D\) is invertible, and \(R_0\) is the direct sum
of the two rectangular maps \(A,F\). In particular the
charge-zero nullity is at least \(M\), with equality if
and only if the two separate rank assertions hold.
When equality holds, the first-order kernel-to-cokernel
map of \(\mathcal J_\lambda\) at zero is precisely
\(\overline C\) under the displayed identifications:
the derivative of the \(\lambda E\) block acts on a zero
\(E_A\otimes Q\) kernel component and contributes nothing.

\paragraph{What has and has not been proved.}
The normal-form identity, (CB.6), and (CB.10)--(CB.15) are
uniform statements for every \(m\). They reduce the endpoint
question to a square global map \(R_1\) without hypotheses.
They reduce it to the \(M\)-dimensional map \(\overline C\)
only after the two separate rank assertions in (CB.14)
have been established. This section does not establish
those assertions or the final isomorphism for arbitrary
\(m\).

For the fixed parameters \(a_i=i\) and \(z=p_m\), the script
\path{agents/jet_review/check_residue_charge_lifting.py}
computes the exact modular kernel and cokernel of the original
residue Jacobian at charge zero, then pairs its charge
derivative with those bases. Its saved certificate
\path{agents/jet_review/residue_charge_lifting_certificate.json}
records, for \(m=3,4,5\), nullities \(3,6,10\) and induced
ranks \(3,6,10\), respectively, modulo \(1000003\).
These verify all three conditions in those sizes only.
They are mechanism checks, not an arbitrary-size proof.

All maps arise from the actual permanent coefficient family
in QA033; their coefficient/comodule interpretation is the
one already fixed there and in F-016. O-031 records the exact
SGA 3, I, Propositions 4.7.2--4.7.3, pp.~34--36, locators.
The block argument itself is finite-dimensional linear algebra
over the indicated characteristic-zero parameter field.
